Municipal 311 systems receive millions of service requests each year ---
heating failures, water leaks, street damage, sanitation problems, noise
--- and agencies plan finite daily service capacity against this demand.
Short-horizon forecasts of request volume are a natural planning input,
and a growing literature shows that public signals such as weather and
calendar structure improve such forecasts. But planners do not consume
mean absolute error; they consume allocations, and the mapping from
predictive accuracy to allocation quality is neither linear nor
guaranteed.

This paper makes exactly three contributions. \textbf{(1)} A reproducible
four-city benchmark of next-day reported 311 demand --- New York, Chicago,
San Francisco, and Austin, 2020--2025 --- built entirely from official
open-data APIs with versioned harmonization rules, provenance manifests,
and SHA-256 checksums, regenerable by one command. \textbf{(2)} A
leakage-controlled evaluation of local, pooled, transfer, calendar,
weather, and probabilistic next-day forecasting, with frozen pre-test
environments, validation-only selection, and temporally censored pooling
and transfer. \textbf{(3)} A controlled, explicitly simulated family-level
capacity-allocation study, under pre-specified hypothetical
service-pressure regimes, that isolates when and how forecast
\emph{representation} changes downstream simulated decision value. The
governing question is:

\begin{quote}
Across heterogeneous municipal 311 reporting systems, under pre-specified
hypothetical service-pressure regimes, when do differences in next-day
forecasts of reported administrative demand produce materially different
outcomes in a transparent simulated family-level capacity allocation, and
when are predictive rankings sufficient proxies for simulated decision
rankings?
\end{quote}

Four design commitments distinguish the evaluation protocol.
\textbf{(i) Frozen environments:} capacity budgets derive from training
data only and are identical at model-selection time and at the single
test evaluation. \textbf{(ii) Validation-only selection:} no headline
configuration is chosen using test outcomes. \textbf{(iii) Temporal
censoring:} pooled and transfer models never train on rows
contemporaneous with any evaluation window. \textbf{(iv) Controlled
uncertainty contrast:} the value of distributional information is
measured within one fitted quantile model --- its median, implied-mean,
and full-distribution representations --- holding everything else fixed.
A guard suite of automated tests enforces each commitment and fails the
build if any is violated.

The contribution is one of data and evaluation rather than of a new
algorithm, and it sits in the big-data tradition along four of the usual
axes. \emph{Volume:} tens of millions of acquired municipal records (36.8
million acquired, 30.2 million retained). \emph{Variety:} four
heterogeneous 311 systems with distinct intake taxonomies, harmonized
service families, calendar structure, and NOAA station weather.
\emph{Veracity:} provenance manifests, documented exclusions, SHA-256
checksums, leakage guards, and one-command reproduction make the assembled
benchmark auditable. \emph{Value:} we measure whether forecasting
improvements do or do not translate into simulated allocation outcomes.
Velocity is deliberately out of scope --- the modeling unit is the daily
aggregate, and we claim no real-time, streaming, or deployed system.

Throughout, we are explicit about what the data can and cannot support.
Reported requests are not social need \citep{kontokosta2021bias}; the
allocation layer is a simulation over abstract request-equivalent
capacity units whose parameters are varied over a pre-specified grid rather than asserted; and no
causal, deployment, or operational-guidance claim is made. The study is
an empirical bridge between the forecasting literature, which stops at
accuracy, and the decision-focused learning literature
\citep{elmachtoub2022smart, mandi2024decision}, which rarely touches
public-sector data.
